\SourcePage{240}{245}
\section{Photoeffect. Relativistic case}

Let us now turn to the case
\begin{equation}
\omega \gg I.
\tag{57.1}
\end{equation}
In this case also $\varepsilon = \omega - I \gg I$, and therefore the influence of the Coulomb field of the nucleus on the wave function of the photoelectron ($\psi'$) can be taken into account with the help of perturbation theory. We write $\psi'$ in the form
\begin{equation}
\psi' = \frac{1}{\sqrt{2\varepsilon}}\,(u'e^{i\mathbf{p}\mathbf{r}} + \psi^{(1)}).
\tag{57.2}
\end{equation}
The photoelectron may be relativistic; therefore the unperturbed function in~(57.2) is written in the form of a relativistic plane wave~(23.1).

Although in the initial state the electron is non-relativistic, its wave function $\psi$ must none the less include (for reasons to be clarified below) a relativistic correction ($\sim Ze^2$). Such a function is given by the formula (see the problem to~\S\,39)
\begin{equation}
\psi = \left(1 - \frac{i}{2m}\,\gamma^0\boldsymbol{\gamma}\nabla\right)\frac{u}{\sqrt{2m}}\,\psi_\mathrm{nr},
\tag{57.3}
\end{equation}
where $\psi_\mathrm{nr}$ is the non-relativistic function of the bound state~(56.6), and $u$ is a bispinor amplitude of the electron at rest, normalised by the condition we have adopted $\bar{u}u = 2m$.

\SourcePage{241}{246}
Substituting the functions~(57.2), (57.3) into the matrix element~(56.2)\,\footnote{Formula~(57.3) was derived for distances $r \sim 1/(mZe^2)$, at which the correction term in it and the relativistic correction to the purely exponential function~(56.6) are always proportional to $Ze^2$. But for the ground state (and indeed for all $s$-states) formula~(57.3) is suitable for any $r$, since the derivative of the purely exponential function~(56.6) (and the correction term in~(57.3)) is always proportional to $Ze^2$.}:
\begin{align}
M_{fi} &= \frac{1}{\sqrt{2m\varepsilon}}\int\Bigl\{\bar{u}'(\boldsymbol{\gamma}\mathbf{e})\left[\left(1 - \frac{i}{2m}\,\gamma^0\boldsymbol{\gamma}\nabla\right)u\psi_\mathrm{nr}\right]e^{-i(\mathbf{p}-\mathbf{k})\mathbf{r}} +{}\notag\\[4pt]
&\qquad\qquad{} + \bar{\psi}^{(1)}_\mathbf{k}(\boldsymbol{\gamma}\mathbf{e})\,e^{i\mathbf{k}\mathbf{r}}\,u\psi_\mathrm{nr}\Bigr\}d^3x.
\tag{57.4}
\end{align}
Having in mind extracting the first term of the expansion of this quantity in powers of $Ze^2$, we can in the second term in the curly brackets replace $\psi_\mathrm{nr}$ simply by the constant $(Ze^2m)^{3/2}/\sqrt{\pi}$. The first term of the result of such a replacement would vanish (for $\mathbf{p} - \mathbf{k} \neq 0$), and it is precisely for this reason that in $\psi$ one must also include the first relativistic correction, proportional to $Ze^2$; for $v \sim 1$ this correction contributes to the cross section terms of the same order as the next term of the expansion of $\psi_\mathrm{nr}$ in powers of $Ze^2$.

In the first term of~(57.4) we carry out the integration by parts, transferring the action of the operator $\nabla$ from $\psi_\mathrm{nr}$ onto the exponential factor. As a result we obtain
\begin{equation}
M_{fi} = \frac{(Ze^2m)^{3/2}}{2(\pi m\varepsilon)^{1/2}}\left\{\bar{u}'(\boldsymbol{\gamma}\mathbf{e})\left[1 + \frac{1}{2m}\,\gamma^0\boldsymbol{\gamma}(\mathbf{p}-\mathbf{k})\right]u\,(e^{-Ze^2mr})_{\mathbf{p}-\mathbf{k}} + \bar{\psi}^{(1)}_{-\mathbf{k}}(\boldsymbol{\gamma}\mathbf{e})u\right\},
\tag{57.5}
\end{equation}
where the vector index denotes the spatial Fourier component. To the accuracy of terms $\sim Ze^2$\,\footnote{Taking the Fourier component from both sides of the equality
\[
(\Delta - \lambda^2)\,\frac{e^{-\lambda r}}{r} = -4\pi\delta(\mathbf{r}),
\]
we obtain
\begin{equation}
\left(\frac{e^{-\lambda r}}{r}\right)_\mathbf{q} = \frac{4\pi}{\mathbf{q}^2 + \lambda^2}.
\tag{57.6a}
\end{equation}
Differentiating this expression with respect to the parameter $\lambda$ gives
\begin{equation}
(e^{-\lambda r})_\mathbf{q} = \frac{8\pi\lambda}{(\mathbf{q}^2 + \lambda^2)^2}.
\tag{57.6b}
\end{equation}}
\begin{equation}
(e^{-Ze^2mr})_{\mathbf{p}-\mathbf{k}} = \frac{8\pi Ze^2m}{(\mathbf{p}-\mathbf{k})^4}.
\tag{57.6}
\end{equation}

\SourcePage{242}{247}
To compute the Fourier component $\bar{\psi}^{(1)}_\mathbf{k}$ we write the equation satisfied by the function $\psi^{(1)}$:
\[
(\gamma^0\varepsilon + i\boldsymbol{\gamma}\nabla - m)\psi^{(1)} = e(\gamma^\mu A_\mu)\,u'e^{i\mathbf{p}\mathbf{r}} = -\frac{Ze^2}{r}\,\gamma^0 u'\,e^{i\mathbf{p}\mathbf{r}}
\]
(this is obtained by substituting~(57.2) into~(32.1)). Applying to both sides of this equation the operator $(\gamma^0\varepsilon + i\boldsymbol{\gamma}\nabla + m)$, we get
\[
(\Delta + \mathbf{p}^2)\psi^{(1)} = -Ze^2(\gamma^0\varepsilon - \boldsymbol{\gamma}\mathbf{k} + m)(\gamma^0 u')\,\frac{1}{r}\,e^{i\mathbf{p}\mathbf{r}}.
\]
Multiplying this equation by $e^{-i\mathbf{k}\mathbf{r}}$ and integrating over $d^3x$, with the terms containing $\Delta$ and $\nabla$ being integrated by parts in the usual way:
\[
(\mathbf{p}^2 - \mathbf{k}^2)\bar{\psi}^{(1)}_\mathbf{k} = -Ze^2(\gamma^0\varepsilon - \boldsymbol{\gamma}\mathbf{k} + m)(\gamma^0 u')\left(\frac{1}{r}\right)_{\mathbf{k}-\mathbf{p}} =
\]
\[
= -Ze^2(2\varepsilon\gamma^0 - \boldsymbol{\gamma}(\mathbf{k}-\mathbf{p}))(\gamma^0 u')\;\frac{4\pi}{(\mathbf{k}-\mathbf{p})^2}.
\]
In the last line we have taken into account that the amplitude $u'$ satisfies the equation
\[
(\varepsilon\gamma^0 - \mathbf{p}\boldsymbol{\gamma} - m)u' = 0,\qquad\text{or}\qquad (\varepsilon\gamma^0 + \mathbf{p}\boldsymbol{\gamma} - m)\gamma^0 u' = 0.
\]
From this we find
\begin{equation}
\bar{\psi}^{(1)}_{-\mathbf{k}}\gamma^0 = 4\pi Ze^2\,\bar{u}'\,\frac{2\varepsilon\gamma^0 + \boldsymbol{\gamma}(\mathbf{k}-\mathbf{p})}{(\mathbf{k}^2 - \mathbf{p}^2)(\mathbf{k}-\mathbf{p})^2}\,\gamma^0.
\tag{57.7}
\end{equation}

Substituting~(57.6), (57.7) into the matrix element~(57.5), we represent it in the form
\[
M_{fi} = \frac{4\pi^{1/2}(Ze^2m)^{5/2}}{(\varepsilon m)^{1/2}(\mathbf{k}-\mathbf{p})^2}\;\bar{u}'Au,
\]
where
\[
A = a(\boldsymbol{\gamma}\mathbf{e}) + (\boldsymbol{\gamma}\mathbf{e})\gamma^0(\boldsymbol{\gamma}\mathbf{b}) + (\boldsymbol{\gamma}\mathbf{c})\gamma^0(\boldsymbol{\gamma}\mathbf{e}),
\]
\[
a = \frac{1}{(\mathbf{k}-\mathbf{p})^2} + \frac{\varepsilon}{m}\,\frac{1}{\mathbf{k}^2 - \mathbf{p}^2},\quad \mathbf{b} = -\frac{\mathbf{k}-\mathbf{p}}{2m(\mathbf{k}-\mathbf{p})^2},\quad \mathbf{c} = \frac{\mathbf{k}-\mathbf{p}}{2m(\mathbf{k}^2 - \mathbf{p})^2}.
\]

The cross section
\[
d\sigma = \frac{8e^2(Ze^2m)^5|\mathbf{p}|}{m\omega(\mathbf{k}-\mathbf{p})^4}\;(\bar{u}'Au)(\bar{u}\bar{A}u')\,do,
\]
where $\bar{A} = \gamma^0 A^+\gamma^0$ (see~\S\,65). This expression must still be summed over the final and averaged over the initial directions of the spin of the electron. These operations are carried out according to the rules described below in~\S\,65, with the help of the polarisation density matrices of the initial and final states:
\[
\rho = \frac{m}{2}\,(\gamma^0 + 1),\qquad \rho' = \frac{1}{2}\,(\gamma^0\varepsilon - \boldsymbol{\gamma}\mathbf{p} + m)
\]

\SourcePage{243}{248}
(in the initial state $\mathbf{p} = 0$, $\varepsilon = m$). They lead to the expression
\[
d\sigma = \frac{16e^2(Ze^2m)^5|\mathbf{p}|}{m\omega(\mathbf{k}-\mathbf{p})^4}\;\mathrm{Sp}\,(\rho'A\rho\bar{A})\,do.
\]

The computation of the trace (using formulae~(22.22)) is a purely algebraic operation and leads to the following result:
\[
\mathrm{Sp}\,(\rho'A\rho\bar{A}) = \frac{m}{\varepsilon+m}\,[a\mathbf{p} - (\mathbf{b} - \mathbf{c})(\varepsilon + m)]^2 + 4m(\mathbf{be})[(\varepsilon + m)(\mathbf{ce}) + a(\mathbf{pe})]
\]
(vector $\mathbf{e}$ is assumed to be real linear polarisation of the photon).

Let us bring the formula for the cross section of the photoeffect to its final form, introducing the polar angle $\theta$ and the azimuth $\varphi$ of the direction $\mathbf{p}$ relative to the direction $\mathbf{k}$, and $\mathbf{e}$ in the plane $xz$ (so that $\mathbf{pe} = |\mathbf{p}|\cos\varphi\sin\theta$). For $\omega \gg I$ the conservation of energy can be written as $\varepsilon - m = \omega$ (instead of $\varepsilon - m = \omega - I$). It is easy to verify that
\[
\mathbf{k}^2 - \mathbf{p}^2 = -2m(\varepsilon - m),\qquad (\mathbf{k}-\mathbf{p})^2 = 2\varepsilon(\varepsilon - m)(1 - v\cos\theta),
\]
where $\mathbf{v} = \mathbf{p}/\varepsilon$ is the velocity of the photoelectron. After straightforward transformations we finally obtain
\begin{align}
d\sigma &= Z^5\alpha^4 r^2_e\;\frac{v^3(1-v^2)^{3/2}\sin^2\theta}{(1-\sqrt{1-v^2}\,)^5(1-v\cos\theta)^4}\times{}\notag\\[4pt]
{}\times{}&\left\{\frac{(1-\sqrt{1-v^2}\,)^2}{2(1-v^2)^{3/2}}\,(1-v\cos\theta) + \left[2 - \frac{(1-\sqrt{1-v^2}\,)(1-v\cos\theta)}{1-v^2}\right]\cos^2\varphi\right\}do,
\tag{57.8}
\end{align}
where $r_e = e^2/m$.

In the ultra-relativistic case ($\varepsilon \gg m$) the cross section of the photoeffect has a sharp maximum at small angles ($\theta \sim \sqrt{1-v^2}$), i.e.\ the electrons are emitted predominantly in the direction of incidence of the photon. Near the maximum we write
\[
1 - v\cos\theta \approx \tfrac{1}{2}[(1 - v^2) + \theta^2],
\]
and the leading terms in~(57.8) give
\begin{equation}
d\sigma \approx 4Z^5\alpha^4 r^2_e\;\frac{(1-v^2)^{3/2}}{(1-v^2+\theta^2)^3}\;\theta\,d\theta\,d\varphi.
\tag{57.9}
\end{equation}

\SourcePage{244}{249}
The elementary, though rather lengthy, integration of expression~(57.8) over angles leads to the following formula for the total cross section (\textit{F.\,Sauter}, 1931):
\begin{align}
\sigma &= 2\pi Z^5\alpha^4 r^2_e\;\frac{(\gamma^2-1)^{3/2}}{(\gamma-1)^5}\times{}\notag\\[4pt]
{}\times{}&\left\{\frac{3}{4} + \frac{\gamma(\gamma-2)}{\gamma+1}\left(1 - \frac{1}{2\gamma\sqrt{\gamma^2-1}}\,\ln\frac{\gamma+\sqrt{\gamma^2-1}}{\gamma-\sqrt{\gamma^2-1}}\right)\right\},
\tag{57.10}
\end{align}
where for brevity we have introduced the ``Lorentz factor''
\begin{equation}
\gamma = \frac{1}{\sqrt{1-v^2}} = \frac{\varepsilon}{m} \approx \frac{m+\omega}{m}.
\tag{57.11}
\end{equation}
In the ultra-relativistic case this formula reduces to the simple expression
\begin{equation}
\sigma = 2\pi Z^5\alpha^4 r^2_e/\gamma.
\tag{57.12}
\end{equation}
In the case $I \ll \omega \ll m$ formula~(57.10) goes to the limit of small $\gamma - 1$ and leads to the result~(56.14), already known to us.
